\setlength\beforechapskip{-1em}
\chapter{Convergence in distribution of Fractional Gaussian Fields}
\chaptermark{Convergence of Fractional Gaussian Fields}

Now that we have presented and proved the existence of fractional Gaussian fields $X_s$ on $K$ and discrete fractional Gaussian fields $X_s^{m}$ on $V_m$, our next task is proving that they converge. Therefore, in this section our main goal is to show that $X_s^{m}$ converges in distribution to $X_s$.

\section{Preliminary Lemmas}
In this section we will prove a collection of lemmas that will be needed.

\begin{lemma}\label{suckItKigami}
	Let $\cbr{P_t}_{t \geq 0}$ be the semigroup generated by the Dirichlet Laplacian $(-\Delta)^{-s}$ and $p_t(x,y)$ the heat kernel admitted by $\cbr{P_t}_{t \geq 0}$. If $\phi_i$ is an orthonormal eigenfunction of $(-\Delta)^{-s}$ then
	\begin{equation}
		P_t \phi_i(x) = e^{-\lambda_i t}\phi_i(x).
	\end{equation}
\end{lemma}
\begin{proof}
	 The lemma follows immediately from these calculations:
	\begin{multline}
		P_t \phi_i(x) = \int_{K}p_t(x,y)\phi_i(y) \dd \mu(y) = \int_{K} \left(\sum_{j = 1}^{\infty}e^{-\lambda_j} \phi_j(x)\phi_j(y)\right) \phi_i(y) \dd \mu(y) \\
		= \int_{K} e^{-\lambda_i t}\phi_i(x) \dd \mu(y) = e^{-\lambda_i t}\phi_i(x).
	\end{multline}
\end{proof}

\begin{lemma}\label{convergenceOfSemigroups}
	For all $f \in S(K)$ and $t > 0$,
	\begin{equation}
		\lim_{m \to \infty} \frac{1}{a_m} \sum_{p \in V_m} f_m(p)P_t^{m} f_m(p) = \int_{K}f(x)P_tf(x) \dd\mu (x)
	\end{equation}
	where $f_m = f \mid_{V_m}$.
\end{lemma}
\begin{proof}
	In this proof we wish to use \cite[Theorem 2.1]{KURTZ1969354}  so we will check that $\Delta$ fulfils the criteria of \cite[Theorem 2.1]{KURTZ1969354}. First there is a bounded linear operator such that for all $f \in L^{2}(K, \mu)$ we have
	$$\lim_{n \to \infty} \norm{P_n f} = \norm{P_n}.$$
	 This holds in our case, since we have $P_n f = f \mid_{V_n} = f_n$.
	
	Secondly we note that $\Delta$ is the infinitesimal generator of a uniformly continuous semigroup by \Cref{genSemigroup}.
	
	Next we have that the sequence of $\cbr{\Delta_m}_{m \geq 0}$ indeed coincides with the extended limit defined, in \cite[355]{KURTZ1969354}, this follows from \Cref{pointwiseFormula} and \Cref{constructDirichletLaplacian}. 
	
	Finally the range $\mathcal{R}(\lambda_0 - \Delta)$ is dense in $L^2$ by \cite[Lemma 1.1.1]{fabriceDirichelt}.
	
	\noi Thus by \cite[Theorem 2.1]{KURTZ1969354}  we have that for all $t \geq 0$
	\begin{equation}\label{convergensofPt}
		\lim_{m \to \infty} \sup_{p \in V_m} \abs{P_t^{m}f_m(p) - \left(P_t f(p)\right)_m} = 0.
	\end{equation}
	If we now write 
	\begin{equation}
		\frac{1}{a_m} \sum_{p \in V_m} f_m(p) P_t^{m}f_m(p) = \int_{V_m} f_m P_t^{m} f_m \dd \mu_m,
	\end{equation}
	then
	\begin{equation}
		\int_{V_m} f_m P_t^{m} f_m \dd\mu_m = \int_{V_m} f_m \left(P_t^{m} f_m - (P_t f)_m \right) \dd\mu_m + \int_{V_m} f_m (P_t f)_m \dd \mu_m .
	\end{equation}
	The first integral converges to $0$ as $m\to \infty$ because of \eqref{convergensofPt}. Furthermore, by \Cref{convergenceOfMeasures} we have that
	\begin{equation}
		\int_{V_m} f_m (P_t f)_m \dd \mu_m \to \int_{K} fP_tf \dd \mu,\quad \text{as } m \to \infty.
	\end{equation}
\end{proof}

\begin{lemma}\label{convergence}
	For all $f \in S(K)$ and $s \geq 0$ the following holds,
	\begin{equation}
		\frac{1}{a_m} \sum_{p \in V_m} f_m(p) (-\Delta_m)^{-2s}f_m(p) \to \int_{K}f(x)(-\Delta)^{-2s}f(x) \dd \mu(x), \quad \text{as } m \to \infty.
	\end{equation}
\end{lemma}
\begin{proof}
	For $s = 0$  we find that 
	\begin{align}
		(-\Delta_m)^{-0} f&= \sum_{i = 1}^{N_m} \frac{1}{(\lambda_i^{m})^{0}} \phi_{i}^{m}(x) \frac{1}{a_m} \sum_{y \in V_m} \phi_{i}^{m}(y) f(y) \\
		&= \sum_{i = 1}^{N_m} \phi_{i}^{m}(x) \int_{V_m} \phi_i^{m}(y)f(y) \dd \mu_{m} \\
		& = f_m.
	\end{align}
	Likewise by \Cref{fracLaplacian} we find
	\begin{align}
		(-\Delta)^{-0} f = \sum_{i = 1}^{\infty} \phi_{i} \int_{K} \phi_i(y) f(y) \dd \mu = f,
	\end{align}
	and the lemma for $s = 0$ follows from \Cref{convergenceOfMeasures}. 
	
		\noi Now assume $s > 0$, then by \Cref{remark:DiscreteLaplacianSemiGroup} we have that
	\begin{align}
		(-\Delta_m)^{s}f_m(x) &= \frac{1}{a_m} \sum_{j = 1}^{N_m} (\lambda_j^{m})^{-s} \phi_j^{m}(x) \sum_{y \in V_m} \phi_j(y) f_m(y) \\
		&= \frac{1}{a_m} \sum_{j = 1}^{N_m} \frac{1}{\Gamma(s)}\int_{0}^{\infty} t^{s-1}e^{-\lambda_j t} \dd t \ \phi_j^{m}(x) \sum_{y \in V_m} \phi_j(y) f_m(y) \\
		&= \frac{1}{\Gamma(s)}\int_{0}^{\infty} t^{s-1} \sum_{j = 1}^{N_m} e^{-\lambda_j t} \phi_j^{m}(x)  \frac{1}{a_m}\sum_{y \in V_m} \phi_j^{m}(y) f_m(y) \dd t \\
		&= \frac{1}{\Gamma(s)} \int_{0}^{\infty} t^{s-1} P_t^{m}f(x) \dd t.
	\end{align}
	 Now by Fubini's theorem we have that
	\begin{align}
		\frac{1}{a_m} \sum_{p \in V_m} f_m(p) (-\Delta_{m})^{-2s}f_m(p)  = \frac{1}{\Gamma(2s)} \int_{0}^{\infty} t^{2s-1} \frac{1}{a_m} \sum_{p \in V_m} f_m(p) P_{t}^{m} f_m(p) \dd t
	\end{align}
	By \Cref{suckItKigami} and \Cref{coro:PtIsLinear} we find that
	\begin{align}
		\frac{1}{a_m} \sum_{p \in V_m} f_m(p) P_{t}^{m} f_m(p) &= \frac{1}{a_m} \sum_{p \in V_m} f_m(p) P_{t}^{m} \sum_{j =1}^{\infty} \inn{f_m}{\phi_i} \phi_{i} \\
		&= \frac{1}{a_m} \sum_{p \in V_m} f_m(p) \sum_{j =1}^{\infty} \inn{f_m}{\phi_j} P_{t}^{m} \phi_{j} \\
		& = \frac{1}{a_m} \sum_{p \in V_m} f_m(p) \sum_{j =1}^{\infty} \inn{f_m}{\phi_j} e^{-\lambda_{j}^{m} t} \phi_{j} \\
		&\leq e^{-\lambda_1^{m}t} \frac{1}{a_m} \sum_{p \in V_m} f_m(p)^2,
	\end{align}
	where we can guarantee the inequality for all $j$ since $0 < \lambda_1 < \lambda_2 <\ldots$. 
	
	Note that 
	\begin{equation}
		\sup_{m} \frac{1}{a_m} \sum_{p \in V_m} f_m(p)^2 < \infty,
	\end{equation}
	 since $f \in S(K) \subset C(K)$ hence $f$ is bounded and $\# V_m < \infty$ for all $m$. We can now by dominated convergence, Fubini's theorem, \Cref{coro:ForRieszKernelAndFractionalLaplacian} and \Cref{convergenceOfSemigroups} conclude that 
	\begin{align}
		\lim_{m \to \infty} \frac{1}{a_m} \sum_{p \in V_m} f_m(p) (-\Delta_{m})^{-2s} f_m(p) &= \frac{1}{\Gamma(2s)} \int_{0}^{\infty} t^{2s-1} \lim_{m \to \infty}\frac{1}{a_m} \sum_{p \in V_m} f_m(p) P_{t}^{m} f_m(p) \dd t \\
		&= \frac{1}{\Gamma(2s)} \int_{0}^{\infty} t^{2s-1} \int_{K} f(x) P_t f(x) \dd\mu(x) \dd t \\
		&= \frac{1}{\Gamma(2s)} \int_{K} f(x) \int_{0}^{\infty} t^{2s-1}P_t f(x) \dd t \dd\mu(x) \\
		&= \int_{K} f(x) (-\Delta)^{-2s}f(x) \dd \mu(x)
	\end{align}
	as desired.
\end{proof}


\section{Main Theorem}
With everything in mind we can now prove the convergence in distribution of discrete fractional Gaussian fields to fractional Gaussian fields.

\medskip
\begin{mainTheorem}
	Let $s \geq 0$, then $X_{s}^{m}$ converges to $X_s$ in distribution in $S'(K)$ when $m \to \infty$.
\end{mainTheorem}

\begin{proof}
	Since $S(K)$ is a nuclear space, by \cite[Corollary~2.4]{biermé2017generalizedrandomfieldslevys} it is enough to prove the convergence of the characteristic functional, e.i that for all $f \in S(K)$
	\begin{equation}
		\e \left[\exp \left(itX_{s}^{m}(f_m) \right)\right] \to \e \left[\exp\left(itX_{s}(f)\right)\right], \quad \text{ as } m \to \infty
	\end{equation} 
	where $f_m = f \mid_{V_m}$. It follows from \Cref{theOneWeNeedInTheEndSomehow} that
	\begin{equation}
		\e \left[\exp\left(itX_{s}^{m}(f_m)\right)\right] = \exp \left(-\frac{t^2}{2} \e \left[ \left(X_{s}^{m}(f_m) \right)^2\right]\right) = \exp \left(-\frac{t^2}{2a_m} \sum_{p \in V_m} f_m(p) (-\Delta_m)^{-2s} f_m(p)\right).
	\end{equation}
	Likewise, via \Cref{fracGaussianField} we find
	\begin{equation}
		\e \left[\exp(i t X_{s}(f))\right] = \exp\left(-\frac{t^2}{2} \e \left[(X_{s}(f))^2\right]\right) = \exp\left(-\frac{t^2}{2}\int_{K} f (-\Delta)^{-2s}f \dd \mu\right)
	\end{equation} 
	where we notice by \cref{coro:ForRieszKernelAndFractionalLaplacian} that
	\begin{align}
		\e \left[(X_{s}(f))^2\right] &= \int_{K} (-\Delta)^{-s}f (-\Delta)^{-s}f \dd \mu \\
		& = \int_{K} \int_{K} f(p)f(q) G_{2s}(p,q) \dd\mu(p) \dd \mu(q) \\
		& = \int_{K} f (-\Delta)^{-2s}f \dd \mu.
	\end{align}
	The proof now follows from \cref{convergence}.
\end{proof}